\section{Experiments} \label{sec:experiments}




\textbf{Datasets}: 
The proposed SPG is evaluated using 5 CL datasets. Their statistics are given in \cref{tab:statisctics_datasets}. Below, we use ``-$n$'' to depict that $n$ tasks are created from each dataset ($n$ takes 10 or 20). 
Classes in the first three datasets are split by task, so each task has a disjoint set of classes. On the other hand, all tasks in the last two datasets have the same set of classes.
We refer to the tasks in the former datasets as ``dissimilar tasks'' in which CF is the essential problem to solve, while we regard the tasks in latter datasets as ``similar tasks'' as the ability of knowledge transfer is more important.





\textbf{(1) CIFAR100-$n$ (C-$n$)}:
CIFAR100~\citep{cifar100} is a dataset that has images of 100 classes. We split it into $n$ tasks so that each task has $100/n$ classes. 
\textbf{(2) TinyImageNet-$n$ (T-$n$)}:
TinyImageNet~\citep{tinyimagenet} is a modified subset of the original ImageNet~\citep{imagenet} dataset, and has 200 classes. Each task contains $200/n$ classes.
\textbf{(3) ImageNet-$100$ (I-$100$)}:
ImageNet~\citep{imagenet} contains 1000 classes of objects. We split it to 100 tasks, each of which has 10 classes, to stress-test systems using a large number of tasks and classes.
\textbf{(4) F-CelebA-$n$ (FC-$n$)}:
Federated CelebA~\citep{celebafemnist} is a dataset of celebrities' face images with several attributes. We use it with binary labels indicating whether he/she in the image is smiling or not. Each task consists of images of one celebrity. 
\textbf{(5) F-EMNIST-$n$ (FE-$n$)}:
Federated EMNIST~\citep{celebafemnist} is a dataset that has 62 classes of hand-written symbols written by different persons. Each task consists of hand-written symbols of one person. 


\begin{table}[t]

\vspace{-3mm}

\caption{Statistics of the CL datasets. $n$ can take $10$ and $20$.
Validation sets are used for early stopping. 
}
\label{tab:statisctics_datasets}
\centering
\resizebox{.95\hsize}{!}{
\begin{tabular}{crrrrr}
\toprule
Dataset & \multicolumn{1}{c}{\#Tasks} & \multicolumn{1}{c}{\#Classes per task} & \multicolumn{1}{c}{\#Train} & \multicolumn{1}{c}{\#Validation} & \multicolumn{1}{c}{\#Test} \\ \midrule
C-$n$ & $n$ & $100/n$ & $45,000$ & $5,000$ & $10,000$ \\
T-$n$ & $n$ & $200/n$ & $90,000$ & $10,000$ & $10,000$ \\
I-100 & $100$ & $10$ & $1,000,000$ & $100,000$ & $50,000$ \\
FC-$n$ & $n$ & $2$ & $400n$ & $40n$ & $80n$ \\
FE-$n$ & $n$ & $62$ & $3100n$ & $310n$ & $620n$ \\
\bottomrule
\end{tabular}
}
\vspace{-5mm}
\end{table}





\textbf{Baselines}:
We use 16 baselines. 11 of them are existing classical and most recent task incremental learning (TIL) systems, \textbf{EWC}~\citep{ewc}, \textbf{A-GEM}~\citep{a-gem},  \textbf{SI}~\citep{si}, \textbf{UCL}~\citep{ucl}, \textbf{TAG}~\citep{tag}, \textbf{PGN}~\citep{pgn}, \textbf{PathNet}~\citep{pathnet}, \textbf{HAT}~\citep{hat}, \textbf{CAT}~\citep{cat}, \textbf{SupSup}~\citep{supsup}, and \textbf{WSN}~\citep{wsn}.
Additionally, three simple methods are used for references: multi-task learning (\textbf{MTL}) that trains all the tasks together,
one task learning (\textbf{ONE}) that learns a separate model/network for each task and thus has no CF or KT,
and naive continual learning (\textbf{NCL})  
that learns each new task without taking any care of previous tasks, i.e., no mechanism to deal with CF. 
Since HAT, our main baseline and perhaps the most effective TIL system, adopts AlexNet~\citep{alexnet} as its backbone, all our experiments are conducted with AlexNet. For other baselines, their original codes are used with switching their backbones to AlexNet for fair comparison. 
Furthermore,
to compare our \textbf{\textit{soft-masking}} with the \textbf{\textit{regularization-based approach}} and our gradient-based importance with Fisher information matrix (FI) based importance in EWC, two more baselines \textbf{EWC-GI} and \textbf{SPG-FI} are created. EWC-GI is EWC with its FI based importance replaced by our gradient-based importance (GI) in \cref{sec:computing_gradients}, i.e., the same penalty/regularization in EWC is applied on our accumulated importance, $\bm{\gamma}_i^{\le t-1}$ in \cref{eq:gamma_i_le_t} when learning task $t$ (no soft-masking). SPG-FI is SPG with our gradient-based importance replaced by FI based importance in EWC. 
The network size of each method is presented in \cref{sec:app_network_size}.



\begin{table*}[t]
\centering
\caption{
Accuracy results in percent.
Best methods in each dataset are emphasized in \textbf{bold}, and second best methods are \underline{underlined}. 
}
\label{tab:results_acc}
\resizebox{.95\hsize}{!}{
\begin{tabular}{crrrrrcrrrrc}
\toprule
& \multicolumn{6}{c}{Dissimilar tasks} & \multicolumn{5}{c}{Similar tasks} \\ \cmidrule(lr){2-7} \cmidrule(lr){8-12}
Model &
\multicolumn{1}{c}{C-10} & \multicolumn{1}{c}{C-20} & \multicolumn{1}{c}{T-10} & \multicolumn{1}{c}{T-20} & \multicolumn{1}{c}{I-100} & \textbf{(Avg.)} &
\multicolumn{1}{c}{FC-10} & \multicolumn{1}{c}{FC-20} & \multicolumn{1}{c}{FE-10} & \multicolumn{1}{c}{FE-20} & \textbf{(Avg.)} \\


\midrule


(MTL) &
$76.4\std{0.3}$ & $78.4\std{0.4}$ & 
$52.7\std{0.3}$ & $59.6\std{1.2}$ &
$64.8\std{0.4}$ &
$66.4$ &
$87.5\std{0.7}$ & $88.3\std{0.2}$ &
$86.2\std{0.7}$ & $87.2\std{2.2}$ &
$87.3$ \\ 


(ONE) &
$66.9\std{3.1}$ & $76.5\std{0.7}$ &
$43.5\std{3.0}$ & $54.5\std{0.9}$ &
$49.3\std{0.4}$ &
$58.1$ &
$74.6\std{2.6}$ & $78.8\std{2.3}$ &
$80.9\std{1.5}$ & $79.7\std{1.4}$ &
$78.5$ \\ 


\hdashline


NCL &
$50.9\std{1.7}$ & $54.2\std{5.3}$ &
$37.2\std{0.9}$ & $41.1\std{1.0}$ & 
$30.6\std{1.2}$ &
$42.8$ &
$84.4\std{1.8}$ & $84.1\std{1.4}$ &
$86.1\std{0.8}$ & $86.4\std{0.3}$ &
$85.3$ \\ 


A-GEM &
$50.8\std{1.0}$ & $56.9\std{7.1}$ & 
$36.2\std{0.6}$ & $41.7\std{1.1}$ & 
$32.1\std{1.1}$ &
$43.5$ &
$83.2\std{3.6}$ & $83.2\std{1.9}$ &
$86.6\std{0.2}$ & $86.9\std{0.3}$ &
$85.0$ \\ 


PGN &
$65.1\std{0.6}$ & $75.5\std{0.4}$ &
$44.0\std{0.8}$ & $53.5\std{0.4}$ &
$45.2\std{0.4}$ &
$56.7$ &
$74.7\std{3.5}$ & $74.7\std{2.8}$ &
$82.5\std{1.0}$ & $82.6\std{0.3}$ &
$78.6$ \\


PathNet &
$69.1\std{0.5}$ & $75.5\std{1.0}$ &
$46.0\std{1.5}$ & $51.9\std{1.6}$ &
$42.0\std{1.5}$ &
$56.9$ &
$79.3\std{1.2}$ & $80.5\std{0.4}$ &
$84.3\std{0.2}$ & $84.5\std{0.4}$ &
$82.1$ \\


HAT &
$62.8\std{0.7}$ & $71.8\std{1.1}$ &
$45.5\std{1.0}$ & $51.7\std{2.1}$ &
$45.3\std{1.9}$ &
$55.4$ &
$79.0\std{3.1}$ & $81.9\std{0.7}$ &
$83.8\std{0.9}$ & $84.6\std{0.8}$ &
$82.3$ \\ 


CAT &
$64.2\std{0.6}$ & $73.9\std{1.1}$ &
$43.7\std{0.6}$ & $50.9\std{0.8}$ & 
\multicolumn{1}{c}{N/A} &
N/A &
$82.9\std{1.3}$ & $82.9\std{3.7}$ &
$82.9\std{1.4}$ & $84.1\std{0.9}$ &
$83.2$ \\ 


SupSup &
$66.2\std{0.2}$ & $75.6\std{0.3}$ &
$44.0\std{0.2}$ & $54.1\std{0.3}$ &
$48.6\std{0.1}$ &
$57.7$ &
$75.0\std{2.3}$ & $78.1\std{1.4}$ &
$80.5\std{0.5}$ & $79.7\std{0.2}$ &
$78.3$ \\


UCL &
$64.8\std{0.8}$ & $74.0\std{0.6}$ &
$45.4\std{0.3}$ & $55.1\std{0.5}$ &
$37.4\std{0.6}$ &
$55.4$ &
$86.2\std{0.5}$ & $86.5\std{0.5}$ &
$85.1\std{0.7}$ & $85.0\std{1.6}$ &
$85.7$ \\


SI &
$62.9\std{0.3}$ & $70.3\std{0.7}$ &
$45.9\std{0.6}$ & $52.6\std{0.7}$ &
$44.1\std{0.2}$ &
$55.2$ &
$86.4\std{0.8}$ & $\underline{86.8}\std{0.3}$ &
$\underline{87.6}\std{0.3}$ &  $\textbf{87.9}\std{0.2}$ &
$\underline{87.2}$ \\


TAG &
$60.6\std{0.7}$ & $68.4\std{0.9}$ &
$43.0\std{0.8}$ & $49.5\std{0.4}$ &
$44.9\std{0.3}$ &
$53.3$ &
$74.3\std{3.7}$ & $77.3\std{2.1}$ &
$84.2\std{0.5}$ & $83.8\std{0.4}$ &
$79.9$ \\


WSN &
$\textbf{69.3}\std{0.2}$ & $\textbf{76.9}\std{0.5}$ &
$47.8\std{0.5}$ & $\underline{57.8}\std{0.5}$ &
$51.9\std{0.4}$ &
$\underline{60.7}$ &
$83.9\std{1.2}$ & $84.1\std{0.7}$ &
$85.5\std{0.2}$ & $86.3\std{0.2}$ &
$84.9$ \\


EWC &
$61.6\std{0.9}$ & $60.7\std{2.7}$ & 
$36.5\std{1.1}$ & $41.5\std{1.3}$ & 
$25.4\std{1.4}$ &
$45.1$ &
$81.2\std{3.0}$ & $86.1\std{0.9}$ &
$86.9\std{0.3}$ & $86.8\std{0.6}$ &
$85.2$ \\ 


\hdashline


EWC-GI &
$63.3\std{1.2}$ & $60.1\std{1.9}$ &
$\underline{48.3}\std{1.0}$ & $48.6\std{1.9}$ &
$\underline{52.7}\std{0.2}$ &
$54.6$ &
$83.4\std{3.0}$ & $84.6\std{1.7}$ &
$86.6\std{1.5}$ & $\underline{87.4}\std{0.9}$ &
$85.5$ \\


SPG-FI &
$60.5\std{0.2}$ & $67.7\std{1.0}$ & 
$43.9\std{0.6}$ & $51.2\std{0.8}$ & 
$48.8\std{0.8}$ & 
$54.4$ &
$\underline{86.7}\std{0.6}$ & $86.2\std{0.8}$ & 
$87.5\std{0.3}$ & $\textbf{87.9}\std{0.2}$ & 
$87.1$ \\


\midrule


\textbf{SPG} &
$\underline{67.7}\std{0.3}$ & $\underline{75.9}\std{1.1}$ & 
$\textbf{48.4}\std{0.3}$ & $\textbf{59.1}\std{0.5}$ & 
$\textbf{58.1}\std{0.5}$ & 
$\textbf{61.8}$ &
$\textbf{87.0}\std{0.9}$ & $\textbf{87.1}\std{0.2}$ & 
$\textbf{87.7}\std{0.2}$ & $\textbf{87.9}\std{0.1}$ & 
$\textbf{87.4}$ \\


\bottomrule
\end{tabular}
}
\vspace{-3mm}
\end{table*}




\textbf{Evaluation Metrics}:
The following three metrics are used. Let $\alpha_i^j$ be the test accuracy of task $i$ task just after a model completes task $j$. 

\textbf{(1)} \textit{Accuracy}:
The average of accuracy for all tasks of a dataset after learning the final task. It is computed by $1/T \sum_t^T \alpha_t^T$, where $T$ is the total number of tasks. \\ 
\textbf{(2)} \textit{Forward transfer}:
This measures how much the learning of previous tasks contributes to the learning of the current task. It is computed by $1/T \sum_t^T \left( \alpha_t^t - \beta_t \right)$, where $\beta_t$ represents the test accuracy of task $t$ in the ONE method, which learns each task separately. \\
\textbf{(3)} \textit{Backward transfer}:
This measures how the learning of the current task affects the performance of the previous tasks. Negative values indicate forgetting. It is computed by $1/(T-1) \sum_t^{T-1} \left( \alpha_t^T - \alpha_t^t \right)$.







\subsection{Training Details} \label{sec:training_details}
The networks are trained with SGD by minimizing the cross-entropy loss except for TAG, which uses the RMSProp optimizer as SGD-based TAG has not been provided by the authors.
The mini-batch size is 64 except MTL that uses 640 for its stability to learn more tasks and classes together.
Hyper-parameters, such as the dropout rate or each method's specific hyper-paramters, are searched based on Tree-structured Parzen Estimator~\citep{tpe}.
With the found best hyper-parameters, the experiments are conducted 5 times with different seeds, and the average accuracy and standard deviation are reported.











\subsection{Results} \label{sec:results}






\cref{tab:results_acc,tab:results_fwt,tab:results_bwt} report the accuracy, forward and backward transfer results, respectively. 
Since CAT takes too much time to train, proportionally to the square of the number of tasks, we are unable to get its results for I-100 (ImageNet with 100 tasks) due to our limited computational resources.








\textbf{Dissimilar Tasks} (C-$n$, T-$n$, I-$100$):
MTL performs the best in all cases, but NCL performs poorly due to serious CF (negative backward transfer) as expected. While PGN, PathNet, HAT, CAT, and SupSup can achieve training with no forgetting (0 backward transfer), on average SPG clearly outperforms all of them as their transfer is limited. Although SPG  underperforms PathNet in C-10, its accuracy is markedly lower in the other settings due to PathNet's capacity problem (see \cref{sec:capacity}).
The backward transfer results in \cref{tab:results_bwt} show that SPG has slight forgetting. However, its positive forward transfer results in \cref{tab:results_fwt} more than make up for the forgetting and give SPG the best final accuracy in most cases. 
As we explained in \cref{sec:related_work}, the regularization-based approach is closely related to our work. However, the representative methods, EWC, SI and UCL, perform poorly due to serious CF (see \cref{tab:results_bwt}) that cannot be compensated by their positive forward transfer (see \cref{tab:results_fwt}). 
Although TAG has almost no CF, its forward transfer is limited, resulting in poorer final accuracy.
While SPG underperforms WSN in C-10 and C-20 slightly, SPG outperforms all baselines in other dissimilar task datasets; especially, in the more realistic and difficult dataset I-100, for which SPG is significantly better, outperforming WSN by 6.2\%.

Comparing our gradient based importance (GI) and Fisher information (FI) matrix based importance, we observe that EWC-GI outperforms EWC except for C-20 (EWC is less than 1\% better), and SPG significantly outperforms SPG-FI in all cases. 
Comparing soft-masking and regularization using the same importance measure, we can see SPG is markedly better than EWC-GI, and SPG-FI is also better than EWC except for C-10 (EWC is only 1\% better).
\textcolor{black}{
These results clearly demonstrate that our gradient based importance (GI) and soft-masking are much more effective than standard \textit{\textbf{regularization methods}}.
}







\textbf{Similar Tasks} (FC-$n$, FE-$n$):
\cref{tab:results_acc} shows that SPG achieves the best in all cases due to its strongest \textit{\textbf{positive forward}} (see \cref{tab:results_fwt}) and \textit{\textbf{positive backward knowledge transfer}} ability (see \cref{tab:results_bwt}).
NCL, A-GEM, UCL, and SI also perform well with positive or very little negative backward transfer since the tasks are similar and CF hardly happens.
TAG underperforms them due to it's limited transfer.
PGN, PathNet, HAT, and WSN
have lower accuracy as their forward transfer is limited (i.e., they just reuse learned parameters in the forward pass) and no positive backward transfer.
SupSup, which does not have any mechanism for knowledge transfer, results in much lower performance.
CAT is slightly better due to its stronger positive forward transfer and no negative backward transfer.

Comparing GI and FI, SPG-FI and EWC-GI perform similarly to SPG and EWC as suppressing updates of important parameters becomes less critical for similar tasks and thus the choice of GI or FI is not important.
Comparing soft-masking and regularization, EWC-GI and EWC are worse than SPG and SPG-FI in all cases, indicating that soft-masking is still more effective as regularization may hinder the learning of new tasks more.




\textbf{Summary}:
SPG markedly outperforms all the baselines. 
When tasks are dissimilar, its positive forward transfer capability overcomes its slight forgetting (negative backward transfer) and achieves the best overall results. It has the strong positive forward transfer even with dissimilar tasks, which has not been realized by the other parameter isolation-based baselines.
When tasks are similar, SPG has both positive forward and backward transfer to achieve the best accuracy results. Keeping most parameters trainable in SPG promotes knowledge transfer. Moreover, we observe that soft-masking (SPG and SPG-FI) is better than regularization (EWC-GI and EWC) and that our gradient-based importance (SPG and EWC-GI) is better than FI (SPG-FI and EWC).









\subsubsection{Capacity Consumption} \label{sec:capacity}
The reason that PGN, PathNet, HAT, CAT, and WSN
have lower accuracy on average than SPG despite the fact that they can learn with no forgetting (see backward transfer in \cref{tab:results_bwt}) is mainly because they suffer from the capacity problem, which is also indirectly reflected in lower forward transfer in \cref{tab:results_fwt}.
Although SupSup does not suffer from this problem, its architecture prevents transfer and gives markedly poorer performances especially in similar tasks (see \cref{tab:results_acc}).
As discussed in \cref{sec:introduction}, since these parameter isolation-based methods freeze a sub-network for each task, as more tasks are learned, the capacity of the network left for learning new knowledge becomes less and less 
except for SupSup,
which leads to poorer performance for later tasks. 
\cref{tab:blocking_ratio} shows the percentage of parameters in the whole network that are completely blocked by HAT and SPG (layer-wise results are given in \cref{sec:app_capacity_layerwise}). 
SPG blocks much fewer parameters than HAT does in all cases. This advantage allows SPG to have more flexibility and capacity to learn while mitigating forgetting, which leads to better performances. 










\cref{fig:results_fwt} plots the forward transfer of the limited datasets due to the space limitations (all results are presented in \cref{sec:app_results_fwt}). 
A positive value in the figures means that a method's forward test result (the test result of the task obtained right after the task is learned) is better than ONE, benefited by the forward knowledge transfer from previous tasks.
We can clearly observe a downward trend of these baselines for dissimilar tasks ((a)-(c)).
We believe that this is due to the capacity problem, i.e., as more tasks are learned, they gradually lose their learning capacity. 
On the other hand, SPG shows a upward trend in all cases, even for dissimilar tasks, thanks to its positive forward transfer, which indicates SPG has a much higher capacity to learn.
For \cref{fig:results_fwt}(d), the difference is not obvious as the tasks are similar and the capacity problem is less serious.
However, SPG still keeps the best forward transfer.



\begin{table}[t]
\centering
\caption{
Forward transfer results in percent. 
}
\label{tab:results_fwt}
\resizebox{\hsize}{!}{
\begin{tabular}{crrrrrrrrrrr}
\toprule
& \multicolumn{6}{c}{Dissimilar tasks} & \multicolumn{5}{c}{Similar tasks} \\ \cmidrule(lr){2-7} \cmidrule(lr){8-12}
Model &
\multicolumn{1}{c}{C-10} & \multicolumn{1}{c}{C-20} & \multicolumn{1}{c}{T-10} & \multicolumn{1}{c}{T-20} & \multicolumn{1}{c}{I-100} & \multicolumn{1}{c}{(Avg.)} &
\multicolumn{1}{c}{FC-10} & \multicolumn{1}{c}{FC-20} & \multicolumn{1}{c}{FE-10} & \multicolumn{1}{c}{FE-20} & \multicolumn{1}{c}{(Avg.)} \\
\midrule


NCL &
$-7.1$ & $-2.7$ &
$-3.2$ & $-4.3$ & 
$-0.7$ &
$-3.6$ &
$+7.9$ & $\underline{+6.2}$ &
$+4.1$ & $+5.9$ &
$+6.0$ \\


A-GEM &
$-3.2$ & $-0.9$ & 
$-3.7$ & $-3.5$ & 
$-0.5$ &
$-2.4$ &
$+8.4$ & $+5.4$ &
$+4.3$ & $+6.4$ &
$+6.1$ \\


PGN &
$-1.8$ & $-1.0$ & 
$+0.5$ & $-1.0$ & 
$-3.9$ &
$-1.4$ &
$+0.2$ & $-4.1$ & 
$+1.6$ & $+2.9$ &
$+0.1$ \\


PathNet &
$+2.2$ & $-0.9$ & 
$+2.5$ & $-2.6$ & 
$-7.1$ &
$-1.2$ &
$+4.7$ & $+1.7$ & 
$+3.4$ & $+4.8$ &
$+3.7$ \\


HAT & 
$-4.0$ & $-4.6$ & 
$+2.0$ & $-2.8$ &
$-3.8$ &
$-2.7$ &
$+4.4$ & $+3.2$ &
$+2.9$ & $+4.9$ & 
$+3.8$ \\


CAT &
$-2.7$ & $-2.5$ &
$+0.1$ & $-3.6$ &
N/A &
N/A &
$+8.4$ & $+4.1$ &
$+2.0$ & $+4.4$ & 
$+4.7$ \\


SupSup &
$-0.6$ & $-0.8$ &
$+0.5$ & $-0.3$ &
$-0.4$ &
$-0.4$ &
$+0.5$ & $-0.8$ &
$-0.4$ & $+0.1$ &
$-0.2$ \\


UCL &
$+3.9$ & $\textbf{+6.1}$ &
$+4.9$ & $+7.8$ &
$+2.7$ &
$+5.1$ &
$+8.3$ & $+5.6$ &
$+3.0$ & $+4.7$ &
$+5.4$ \\


SI &
$\textbf{+7.5}$ & $+3.5$ &
$\textbf{+13.3}$ & $\textbf{+8.6}$ &
$\textbf{+12.6}$ &
$\textbf{+9.1}$ &
$+7.1$ & $+5.8$ &
$\textbf{+5.9}$ & $\underline{+7.5}$ &
$+6.6$ \\


TAG &
$-5.6$ & $-6.6$ &
$-0.2$ & $-4.2$ &
$-5.1$ &
$-4.3$ &
$-0.5$ & $-1.8$ &
$+3.5$ & $+4.1$ &
$+1.3$ \\


WSN & 
$+2.4$ & $+0.5$ &
$+4.3$ & $+3.3$ &
$+2.9$ & 
$+2.7$ & 
$+9.3$ & $+5.3$ &
$+4.6$ & $+6.6$ &
$+6.4$ \\

EWC &
$+0.4$ & $-2.5$ & 
$-3.5$ & $-5.6$ &
$-1.2$ &
$-2.5$ &
$+7.8$ & $+5.8$ &
$+4.7$ & $+6.7$ &
$+6.3$ \\


\hdashline


EWC-GI &
$-1.5$ & $+2.3$ &
$+6.3$ & $+6.8$ &
$\underline{+12.1}$ &
$+5.2$ &
$+8.8$ & $+5.7$ &
$+4.2$ & $+6.0$ &
$+6.2$ \\


SPG-FI &
$+3.2$ & $+1.0$ & 
$+1.5$ & $+4.4$ & 
$+6.6$ & 
$+3.4$ &
$\underline{+9.6}$ & $\textbf{+6.5}$ & 
$\underline{+5.4}$ & $+7.4$ & 
$\underline{+7.2}$ \\


\midrule


\textbf{SPG} &
$\underline{+5.5}$ & $\underline{+5.0}$ & 
$\underline{+8.9}$ & $\underline{+7.9}$ & 
$+10.2$ & 
$\underline{+7.5}$ &
$\textbf{+9.8}$ & $+5.6$ & 
$\textbf{+5.9}$ & $\textbf{+7.8}$ & 
$\textbf{+7.3}$ \\

\bottomrule
\end{tabular}
}
\vspace{-5mm}
\end{table}




\begin{table}[t]
\centering
\caption{
Backward transfer results in percent. 
}
\label{tab:results_bwt}
\resizebox{\hsize}{!}{
\begin{tabular}{crrrrrrrrrrr}
\toprule
& \multicolumn{6}{c}{Dissimilar tasks} & \multicolumn{5}{c}{Similar tasks} \\ \cmidrule(lr){2-7} \cmidrule(lr){8-12}
Model & 
\multicolumn{1}{c}{C-10} & \multicolumn{1}{c}{C-20} & \multicolumn{1}{c}{T-10} & \multicolumn{1}{c}{T-20} & \multicolumn{1}{c}{I-100} & \multicolumn{1}{c}{(Avg.)} &
\multicolumn{1}{c}{FC-10} & \multicolumn{1}{c}{FC-20} & \multicolumn{1}{c}{FE-10} & \multicolumn{1}{c}{FE-20} & \multicolumn{1}{c}{(Avg.)} \\
\midrule


NCL &
$-9.9$ & $-20.5$ &
$-3.4$ & $-9.6$ & 
$-17.9$ &
$-12.3$ &
$+2.1$ & $-0.9$ &
$+1.2$ & $\underline{+0.9}$ &
$+0.8$ \\


A-GEM &
$-14.3$ & $-21.8$ & 
$-4.0$ & $-9.7$ & 
$-16.6$ &
$-13.3$ &
$+0.2$ & $-1.0$ &
$\underline{+1.5}$ & $\underline{+0.9}$ & 
$+0.4$ \\


PGN &
$\textbf{0.0}$ & $\textbf{0.0}$ & 
$\textbf{0.0}$ & $\textbf{0.0}$ & 
$\underline{0.0}$ & 
$\textbf{0.0}$ &
$0.0$ & $0.0$ & 
$0.0$ & $0.0$ &
$0.0$ \\


PathNet &
$\textbf{0.0}$ & $\textbf{0.0}$ & 
$\textbf{0.0}$ & $\textbf{0.0}$ & 
$\underline{0.0}$ &
$\textbf{0.0}$ &
$0.0$ & $0.0$ &
$0.0$ & $0.0$ &
$0.0$ \\


HAT &
$\textbf{0.0}$ & $\textbf{0.0}$ &
$\textbf{0.0}$ & $\textbf{0.0}$ &
$\underline{0.0}$ &
$\textbf{0.0}$ &
$0.0$ & $0.0$ &
$0.0$ & $0.0$ & 
$0.0$ \\


CAT &
$\textbf{0.0}$ & $\textbf{0.0}$ & 
$\textbf{0.0}$ & $\textbf{0.0}$ &
N/A &
N/A &
$0.0$ & $0.0$ &
$0.0$ & $0.0$ & 
$0.0$ \\


SupSup &
$\textbf{0.0}$ & $\textbf{0.0}$ & 
$\textbf{0.0}$ & $\textbf{0.0}$ &
$\underline{0.0}$ & 
$\textbf{0.0}$ &
$0.0$ & $0.0$ &
$0.0$ & $0.0$ &
$0.0$ \\


UCL &
$-6.6$ & $-9.0$ &
$-3.4$ & $-7.7$ &
$-14.5$ &
$-8.2$ &
$\underline{+3.7}$ & $+2.2$ &
$+1.2$ & $+0.7$ & 
$\underline{+2.0}$ \\


SI &
$-12.8$ & $-5.7$ &
$-12.1$ & $-11.0$ &
$-17.6$ &
$-11.8$ &
$\textbf{+5.2}$ & $\underline{+2.3}$ &
$+0.9$ & $+0.7$ & 
$\textbf{+2.3}$ \\


TAG &
$\underline{-0.8}$ & $\underline{-1.6}$ &
$\underline{-0.4}$ & $\underline{-0.8}$ &
$\textbf{+0.9}$ &
$\underline{-0.5}$ &
$+0.4$ & $+0.4$ &
$-0.3$ & $0.0$ & 
$+0.1$ \\


WSN & 
$\textbf{0.0}$ & $\textbf{0.0}$ &
$\textbf{0.0}$ & $\textbf{0.0}$ &
$\underline{0.0}$ & 
$\textbf{0.0}$ & 
$0.0$ & $0.0$ &
$0.0$ & $0.0$ &
$0.0$ \\


EWC & 
$-6.4$ & $-13.9$ & 
$-3.9$ & $-7.8$ & 
$-22.7$ &
$-10.9$ &
$-1.3$ & $+1.6$ &
$+1.4$ & $+0.4$ & 
$+0.5$ \\


\hdashline


EWC-GI &
$-1.4$ & $-19.6$ &
$-1.6$ & $-13.4$ &
$-8.6$ &
$-8.9$ &
$0.0$ & $+0.1$ &
$\textbf{+1.6}$ & $\textbf{+1.7}$ &
$+0.9$ \\


SPG-FI &
$-10.7$ & $-10.3$ &
$-1.3$ & $-8.1$ &
$-7.0$ & 
$-7.5$ &
$+2.8$ & $+1.0$ &
$+1.4$ & $+0.8$ &
$+1.5$ \\


\midrule


\textbf{SPG} &
$-5.3$ & $-5.9$ &
$-4.4$ & $-3.4$ &
$-1.2$ &
$-4.0$ &
$+2.9$ & $\textbf{+2.8}$ &
$+0.9$ & $+0.5$ &
$+1.8$ \\


\bottomrule
\end{tabular}
}
\vspace{-3mm}
\end{table}



\begin{table}[h]
\centering
\caption{
Each cell reports how many percentage of parameters are completely blocked (i.e., importance of 1) just after learning task $t$. $T$ is the total number of tasks (e.g., $T=10$ for C-10).
}
\label{tab:blocking_ratio}
\resizebox{\columnwidth}{!}{
\begin{tabular}{ccrrrrrrrrr}
\toprule
& & \multicolumn{5}{c}{Dissimilar tasks} & \multicolumn{4}{c}{Similar tasks} \\ \cmidrule(lr){3-7} \cmidrule(lr){8-11}
$t$ & Model &
\multicolumn{1}{c}{C-10} & \multicolumn{1}{c}{C-20} & \multicolumn{1}{c}{T-10} & \multicolumn{1}{c}{T-20} & \multicolumn{1}{c}{I-100} & \multicolumn{1}{c}{FC-10} & \multicolumn{1}{c}{FC-20} & \multicolumn{1}{c}{FE-10} & \multicolumn{1}{c}{FE-20}  \\
\midrule

\multirow{2}{*}{$1$}
& HAT & 
$1.9\std{0.5}$ & $15.9\std{1.1}$ & 
$2.9\std{0.2}$ & $0.2\std{0.1}$ &
$23.5\std{0.6}$ & 
$0.3\std{0.0}$ & $6.3\std{1.0}$ &
$23.9\std{1.0}$ & $19.0\std{1.1}$ \\
& \textbf{SPG} & 
$0.1\std{0.0}$ & $0.2\std{0.0}$ &
$0.1\std{0.0}$ & $0.1\std{0.0}$ &
$0.1\std{0.0}$ &
$0.1\std{0.0}$ & $0.2\std{0.0}$ &
$0.1\std{0.0}$ & $0.1\std{0.0}$ \\

\hdashline

\multirow{2}{*}{$T/2$}
& HAT & 
$22.4\std{1.1}$ & $98.2\std{0.3}$ &
$36.7\std{1.9}$ & $21.1\std{1.4}$ &
$99.8\std{0.0}$ &
$3.4\std{0.3}$ & $65.7\std{2.4}$ &
$86.8\std{0.8}$ & $94.0\std{1.6}$ \\
& \textbf{SPG} & 
$0.9\std{0.1}$ & $2.9\std{0.3}$ &
$0.5\std{0.1}$ & $1.4\std{0.2}$ &
$4.1\std{0.4}$ &
$0.8\std{0.0}$ & $1.3\std{0.1}$ &
$0.6\std{0.3}$ & $1.0\std{0.2}$ \\

\hdashline

\multirow{2}{*}{$T$}
& HAT & 
$41.8\std{1.3}$ & $99.6\std{0.1}$ &
$57.5\std{1.8}$ & $38.6\std{2.1}$ &
$99.9\std{0.1}$ & 
$9.8\std{1.1}$ & $79.5\std{0.9}$ &
$97.9\std{0.3}$ & $97.5\std{0.9}$ \\
& \textbf{SPG} &
$2.0\std{0.2}$ & $5.2\std{0.6}$ &
$0.9\std{0.1}$ & $3.1\std{0.2}$ &
$6.3\std{0.6}$ &
$1.4\std{0.1}$ & $2.0\std{0.3}$ &
$1.1\std{0.5}$ & $1.5\std{0.3}$ \\

\bottomrule
\end{tabular}
}
\vspace{-2mm}
\end{table}

\begin{figure}[t]

\begin{minipage}{.48\hsize}
\centering
\includegraphics[width=\hsize]{images/fwt_C-20_leg.pdf}
\subcaption{C-20}
\end{minipage}
\hfill
\begin{minipage}{.48\hsize}
\centering
\includegraphics[width=\hsize]{images/fwt_T-20.pdf}
\subcaption{T-20}
\end{minipage}

\medskip

\begin{minipage}{.48\hsize}
\centering
\includegraphics[width=\hsize]{images/fwt_I-100.pdf}
\subcaption{I-100}
\end{minipage}
\hfill
\begin{minipage}{.48\hsize}
\centering
\includegraphics[width=\hsize]{images/fwt_FE-20.pdf}
\subcaption{FE-20}
\end{minipage}

\caption{
The forward transfer of each task along with the number of tasks learned. 
(a) to (c) are the plots with dissimilar tasks, while (d) is the one with similar tasks.
}
\label{fig:results_fwt}

\vspace{-3mm}

\end{figure}












\begin{table*}[t]
\centering
\caption{
The results for pruning parameters based on the gradient-based importance.
}
\label{tab:app_pruning}
\resizebox{.95\hsize}{!}{
 \begin{tabular}{rrrrrrrrrr}
\toprule
& \multicolumn{5}{c}{Dissimilar tasks} & \multicolumn{4}{c}{Similar tasks} \\ \cmidrule(lr){2-6} \cmidrule(lr){7-10}
\multicolumn{1}{c}{Pruning strategy} & \multicolumn{1}{c}{C-10} & \multicolumn{1}{c}{C-20} & \multicolumn{1}{c}{T-10} & \multicolumn{1}{c}{T-20} & \multicolumn{1}{c}{I-100} &
 \multicolumn{1}{c}{FC-10} & \multicolumn{1}{c}{FC-20} & \multicolumn{1}{c}{FE-10} & \multicolumn{1}{c}{FE-20} \\
\midrule

Nothing & 
$75.5\std{1.0}$ & $78.9\std{1.9}$ &
$46.7\std{0.8}$ & $49.2\std{1.2}$ &
$48.8\std{1.2}$ & 
$74.5\std{3.0}$ & $87.8\std{2.6}$ &
$86.0\std{0.5}$ & $85.4\std{0.8}$ \\

\hdashline

Lowest $10$\% &
$73.6\std{2.4}$ & $75.4\std{4.8}$ &
$41.2\std{0.9}$ & $45.8\std{1.5}$ &
$43.3\std{2.6}$ & 
$74.0\std{2.9}$ & $85.5\std{1.9}$ &
$83.1\std{1.7}$ & $84.3\std{0.4}$ \\

Random $10$\% &
$69.5\std{2.1}$ & $70.5\std{4.2}$ &
$34.7\std{1.1}$ & $44.8\std{1.6}$ &
$40.8\std{1.2}$ & 
$71.3\std{2.0}$ & $84.3\std{3.7}$ &
$70.8\std{4.6}$ & $72.3\std{3.1}$ \\

Highest $10$\% &
$19.4\std{8.5}$ & $24.6\std{8.1}$ &
$6.2\std{1.1}$ & $13.2\std{2.4}$ &
$17.1\std{3.6}$ & 
$39.0\std{0.6}$ & $49.3\std{4.1}$ &
$55.5\std{9.7}$ & $19.8\std{8.4}$ \\

\hdashline

Lowest $20$\% &
$68.6\std{5.1}$ & $72.7\std{6.7}$ &
$36.4\std{3.3}$ & $43.1\std{1.8}$ &
$39.4\std{4.1}$ & 
$73.3\std{3.9}$ & $85.5\std{3.7}$ &
$61.5\std{2.7}$ & $81.6\std{3.5}$ \\

Random $20$\% &
$58.3\std{2.2}$ & $58.7\std{7.4}$ &
$20.4\std{0.6}$ & $37.2\std{3.1}$ &
$32.3\std{2.6}$ & 
$68.5\std{3.0}$ & $85.0\std{3.6}$ &
$37.2\std{3.4}$ & $43.5\std{5.4}$ \\

Highest $20$\% &
$10.8\std{1.0}$ & $22.5\std{5.6}$ &
$5.5\std{0.6}$ & $10.2\std{0.2}$ &
$11.7\std{1.3}$ & 
$38.8\std{0.0}$ & $49.3\std{4.1}$ &
$20.7\std{5.8}$ & $7.8\std{1.5}$ \\

\bottomrule
\end{tabular}
}

\vspace{-3mm}

\end{table*}


\subsubsection{Validity of Gradient-based Importance}


We further analyze how the proposed gradient-based importance metric
indicates the contribution of parameters with different importance values to the performance of each task. Specifically, we evaluate the accuracy on the first task of each dataset (no continual learning) 
after pruning some parameters based on their importance so that we can confirm whether the importance metric is co-related with the performance change.

The following four strategies for choosing which parameters to prune are compared --
\textbf{(1) Nothing}: 
we do not prune any parameters. 
\textbf{(2) Random $\bm{n}$\%}: 
$n\%$ of parameters are randomly pruned regardless of their importance.
\textbf{(3) Lowest $\bm{n}$\%}: 
the parameters with the lowest $n\%$ of importance are pruned.
\textbf{(4) Highest $\bm{n}$\%}:
the parameters with the highest $n\%$ of importance are pruned.


The average results over 5 different seeds are presented in \cref{tab:app_pruning}.
It can be clearly observed that pruning parameters with higher importance degrades the performance more (e.g., ``Lowest $10$\%'' is much better than ``Highest $10$\%''). When it is ``Highest 20\%'', the accuracy is almost like random chance classification (e.g., it is 10.8\% for C-10 while the random chance also gives 10\%). We believe this observation implies that our gradient-based importance metric effectively indicates the importance of parameters.


\begin{figure}[t]

\begin{minipage}{.48\hsize}
\centering
\includegraphics[width=\hsize]{images/repr_c10_t_leg.pdf}
\subcaption{Fine-tuning for TinyImageNet after CL for C-10}
\end{minipage}
\hfill
\begin{minipage}{.48\hsize}
\centering
\includegraphics[width=\hsize]{images/repr_i100_c.pdf}
\subcaption{Fine-tuning for CIFAR100 after CL for I-100}
\end{minipage}

\caption{
The learning of representation through continual learning.
The x-axis means the number of tasks learned in continual learning (CL).
The pair of a CL/non-CL dataset for (a) and (b) is C-10/TinyImageNet and I-100/CIFAR100, respectively.
}
\label{fig:cl_representation}

\vspace{-5mm}

\end{figure}



\begin{table*}[tb]
\centering
\caption{
The results for the ablation studies.
}
\label{tab:results_ablation}
\resizebox{.95\hsize}{!}{
\begin{tabular}{lrrrrrcrrrrc}
\toprule
& \multicolumn{6}{c}{Dissimilar tasks} & \multicolumn{5}{c}{Similar tasks} \\ \cmidrule(lr){2-7} \cmidrule(lr){8-12}
\multicolumn{1}{c}{Ablation} & \multicolumn{1}{c}{C-10} & \multicolumn{1}{c}{C-20} & \multicolumn{1}{c}{T-10} & \multicolumn{1}{c}{T-20} & \multicolumn{1}{c}{I-100} & \textbf{(Avg.)} & \multicolumn{1}{c}{FC-10} & \multicolumn{1}{c}{FC-20} & \multicolumn{1}{c}{FE-10} & \multicolumn{1}{c}{FE-20} & \textbf{(Avg.)} \\
\midrule

\textbf{SPG} & 
$\textbf{67.7}\std{0.3}$ & $\textbf{75.9}\std{1.1}$ &
$\textbf{48.4}\std{0.3}$ & $\textbf{59.1}\std{0.5}$ &
$\textbf{58.1}\std{0.5}$ & 
$\textbf{61.8}$ & 
$\textbf{87.0}\std{0.9}$ & $\textbf{87.1}\std{0.2}$ &
$\underline{87.7}\std{0.2}$ & $87.9\std{0.1}$ &
$\textbf{87.4}$ \\

\hdashline

SPG (w/o CHI) & 
$65.4\std{0.6}$ & $70.4\std{2.2}$ &
$46.9\std{0.7}$ & $55.3\std{0.8}$ &
$\underline{57.6}\std{0.4}$ & 
$59.1$ & 
$\underline{86.2}\std{0.3}$ & $\underline{86.9}\std{0.5}$ &
$\textbf{87.9}\std{0.3}$ & $\textbf{88.1}\std{0.3}$ &
$\underline{87.3}$ \\

SPG (w/o SMH) &
$\underline{66.1}\std{1.0}$ & $\underline{74.8}\std{0.3}$ &
$\underline{47.2}\std{0.5}$ & $\underline{56.9}\std{0.4}$ &
$55.2\std{0.8}$ & 
$\underline{60.0}$ & 
$85.8\std{0.7}$ & $86.8\std{0.5}$ &
$87.2\std{0.2}$ & $87.5\std{0.2}$ &
$86.8$ \\

SPG (w/ hard-masking) &
$63.6\std{0.4}$ & $73.2\std{1.2}$ &
$46.8\std{0.4}$ & $53.7\std{0.2}$ &
$51.1\std{0.3}$ &
$57.7$ & 
$84.3\std{0.4}$ & $86.3\std{0.5}$ &
$86.8\std{0.0}$ & $\underline{88.0}\std{0.2}$ &
$86.3$ \\

\bottomrule
\end{tabular}
}

\vspace{-3mm}

\end{table*}



\subsubsection{Better Representation Learning of SPG}

We found that SPG's stronger performance is manifested in its strong representations learning. 
We conduct additional experiments from the perspective of  representation learning.
In particular, a model that has just incrementally learned some tasks of a CL dataset (e.g., 5 tasks of C-10) are used as a frozen feature extractor to learn another dataset (not split into tasks), which we call ``non-CL dataset'', by fine-tuning a new classifier using the dataset, e.g., CIFAR100 or TinyImageNet (a non-CL dataset contains all classes of the dataset shown in \cref{tab:statisctics_datasets}). 
We evaluate the test accuracy for a non-CL dataset. Three pairs of a CL/non-CL dataset are tested: (1) C-10/TinyImageNet, (2) I-100/CIFAR100, and (3) T-10/CIFAR100. We here only show the results for (1) and (2) due to the space limitations (all the results are provided in \cref{sec:app_representation}).


\cref{fig:cl_representation}(a) shows that SPG learns better representations in continual learning than baselines (we use NCL and two strong performing baselines from \cref{tab:results_acc}). HAT even deteriorates, which indicates that hard-masking of some parameters/units leave less network capacity to learn good features. 
\cref{fig:cl_representation}(b) shows that while NCL significantly degrades due to its serious forgetting in such an extreme continual learning on I-100 (with 100 tasks), SPG keeps the best after learning a few tasks.
These clearly confirm that SPG learns better representations to enable better continual learning. 










\subsection{Ablation Studies} \label{sec:ablation}


As SPG has three core mechanisms that contribute to its performance, we evaluate whether each of them is beneficial.
The first one is \textit{\underline{c}ross-\underline{h}ead \underline{i}mportance} (CHI), which is introduced in \cref{eq:gamma_i_t} with the motivation of suppressing the update of important parameters for previous tasks more. In the ablation \textbf{SPG (w/o CHI)}, $\bm{\gamma}_i^t$ is replaced with $\bm{\gamma}_i^{t,t}$ in \cref{eq:gamma_i_t} and the previous tasks' heads are not used for computing importance.
The second one is \underline{\textit{s}}oft-\underline{\textit{m}}asking of each \underline{\textit{h}}ead (SMH), which is introduced in \cref{eq:g_ht_prime} for the purpose of balancing the training of the feature extractor and the classification head. In the ablation \textbf{SPG (w/o SMH)}, \cref{eq:g_ht_prime} is not employed.
The last one is soft-masking (not hard-masking), which is the core technique of SPG to keep most parameters trainable while mitigating CF. While the reported results so far are based on soft-masking, it is also possible to convert the importance to binary (0 or 1) masks using a pre-defined threshold. In the ablation \textbf{SPG (w/ hard-masking)}, if the threshold is 0.6, this variant treats importance larger than 0.6 as 1 (blocking updates of parameters), otherwise 0 (not blocking). We search for the best threshold for each dataset from $\{0.2,0.4,0.6,0.8\}$. 



The ablation results are presented in \cref{tab:results_ablation}.
For CHI, it improves the performance especially in dissimilar tasks (up to 5.5\% for C-20) by blocking more gradient flow to mitigate CF. On the other hand, CHI does not contributes much for similar tasks, which is reasonable as blocking parameters becomes less important in similar tasks. Nevertheless, 
SPG with CHI still works the best for similar tasks on average and CHI does not cause side effects. More quantitative analysis of how CHI contributes to suppressing parameter updates is given in \cref{sec:app_chi}.
For SMH, it delivers positive performance gains both in dissimilar and similar tasks. These results indicate that the lack of balance in the training between the feature extractor and head, which can happen without SMH, adversely affects the performance. The promotion of knowledge transfer by SMH becomes most prominent for I-100 (up to 2.9\%) as the effectiveness of knowledge transfer is more important in such an extreme case with more tasks.
For whether masking should be soft or hard, SPG with hard-masking is significantly worse than SPG with soft-masking, which demonstrates that it is difficult to find a good threshold for hard-masking. Soft-masking is more flexible and effective.
These results clearly confirm that SPG enjoys all of the three different mechanisms.
